\documentclass[12pt,a4paper]{article}
\usepackage{iftex}
\ifPDFTeX
    \usepackage[utf8]{inputenc}
    \usepackage[T2A]{fontenc}
    \usepackage[russian]{babel}
\else
    \usepackage{fontspec}
    \setmainfont{Liberation Serif}
    \setsansfont{Liberation Sans}
    \setmonofont{DejaVu Sans Mono}[Scale=0.9]
    \usepackage[russian,provide=*]{babel}
\fi
\usepackage{amsmath, amssymb}
\usepackage{graphicx}
\usepackage{hyperref}
\usepackage{listings}
\usepackage{xcolor}
\usepackage{geometry}
\usepackage{float}
\usepackage{caption}
\usepackage{enumitem}
\usepackage{titlesec}
\usepackage{fancyhdr}
\usepackage{array}
\usepackage{longtable}
\newcolumntype{L}[1]{>{\raggedright\arraybackslash}p{#1}}
\usepackage{tikz}
\usetikzlibrary{positioning, arrows.meta, fit}

% Настройка полей
\geometry{a4paper, left=2.5cm, right=1.5cm, top=2cm, bottom=2cm}

% Настройка заголовков
\titleformat{\section}{\Large\bfseries}{\thesection}{1em}{}
\titleformat{\subsection}{\large\bfseries}{\thesubsection}{1em}{}

\hypersetup{colorlinks=true, linkcolor=black, urlcolor=blue}
\captionsetup{labelsep=endash}

% Настройка кода (literate нужен, чтобы listings печатал кириллицу в pdfLaTeX)
\lstset{
    language=Python,
    basicstyle=\ttfamily\footnotesize,
    keywordstyle=\color{blue},
    commentstyle=\color{green!40!black},
    stringstyle=\color{red},
    numbers=left,
    numberstyle=\tiny,
    breaklines=true,
    frame=single,
    captionpos=b,
    showstringspaces=false,
    tabsize=4,
    extendedchars=true,
    keepspaces=true,
    columns=flexible,
    breakatwhitespace=true
}
\ifPDFTeX
    \lstset{literate={а}{{\cyra}}1 {б}{{\cyrb}}1 {в}{{\cyrv}}1 {г}{{\cyrg}}1 {д}{{\cyrd}}1 {е}{{\cyre}}1 {ё}{{\cyryo}}1 {ж}{{\cyrzh}}1 {з}{{\cyrz}}1 {и}{{\cyri}}1 {й}{{\cyrishrt}}1 {к}{{\cyrk}}1 {л}{{\cyrl}}1 {м}{{\cyrm}}1 {н}{{\cyrn}}1 {о}{{\cyro}}1 {п}{{\cyrp}}1 {р}{{\cyrr}}1 {с}{{\cyrs}}1 {т}{{\cyrt}}1 {у}{{\cyru}}1 {ф}{{\cyrf}}1 {х}{{\cyrh}}1 {ц}{{\cyrc}}1 {ч}{{\cyrch}}1 {ш}{{\cyrsh}}1 {щ}{{\cyrshch}}1 {ъ}{{\cyrhrdsn}}1 {ы}{{\cyrery}}1 {ь}{{\cyrsftsn}}1 {э}{{\cyrerev}}1 {ю}{{\cyryu}}1 {я}{{\cyrya}}1 {А}{{\CYRA}}1 {Б}{{\CYRB}}1 {В}{{\CYRV}}1 {Г}{{\CYRG}}1 {Д}{{\CYRD}}1 {Е}{{\CYRE}}1 {Ё}{{\CYRYO}}1 {Ж}{{\CYRZH}}1 {З}{{\CYRZ}}1 {И}{{\CYRI}}1 {Й}{{\CYRISHRT}}1 {К}{{\CYRK}}1 {Л}{{\CYRL}}1 {М}{{\CYRM}}1 {Н}{{\CYRN}}1 {О}{{\CYRO}}1 {П}{{\CYRP}}1 {Р}{{\CYRR}}1 {С}{{\CYRS}}1 {Т}{{\CYRT}}1 {У}{{\CYRU}}1 {Ф}{{\CYRF}}1 {Х}{{\CYRH}}1 {Ц}{{\CYRC}}1 {Ч}{{\CYRCH}}1 {Ш}{{\CYRSH}}1 {Щ}{{\CYRSHCH}}1 {Ъ}{{\CYRHRDSN}}1 {Ы}{{\CYRERY}}1 {Ь}{{\CYRSFTSN}}1 {Э}{{\CYREREV}}1 {Ю}{{\CYRYU}}1 {Я}{{\CYRYA}}1 {—}{{\textemdash}}1 {–}{{\textendash}}1 {°}{{\textdegree}}1 {→}{{$\rightarrow$}}1}
\fi
\renewcommand{\lstlistingname}{Листинг}
\lstdefinelanguage{json}{basicstyle=\ttfamily\small, string=[s]{"}{"}, stringstyle=\color{red}, morecomment=[l]{//}}

% Настройка колонтитулов
\pagestyle{fancy}
\fancyhf{}
\fancyhead[L]{Лабораторная работа №3: RAG-система}
\fancyhead[R]{\thepage}

% Скриншот: если файла нет, вместо него рамка-заглушка (документ всё равно собирается)
\newcommand{\screenshot}[3][\textwidth]{%
    \begin{figure}[H]
        \centering
        \IfFileExists{images/#2}{\includegraphics[width=#1]{images/#2}}%
            {\fbox{\parbox[c][3cm][c]{0.9\textwidth}{\centering Место для скриншота: \texttt{images/\detokenize{#2}}}}}
        \caption{#3}
        \label{fig:#2}
    \end{figure}}


\begin{document}
\sloppy
\begin{titlepage}
\begin{center}
\textbf{МИНИСТЕРСТВО НАУКИ И ВЫСШЕГО ОБРАЗОВАНИЯ РОССИЙСКОЙ ФЕДЕРАЦИИ}
\vspace{0.5cm}

\textbf{Федерального государственного автономного образовательного учреждения высшего образования}
\vspace{0.3cm}

\textbf{«НАЦИОНАЛЬНЫЙ ИССЛЕДОВАТЕЛЬСКИЙ ТЕХНОЛОГИЧЕСКИЙ УНИВЕРСИТЕТ МИСИС»}
\vspace{2cm}

\textbf{Институт компьютерных наук}
\vspace{0.3cm}

\textbf{Кафедра Инженерной кибернетики}
\vspace{1cm}

\textbf{Отчёт по лабораторной работе №3}
\vspace{0.3cm}

\textbf{по дисциплине «Современные интеллектуальные сетевые сервисы»}
\vspace{1.5cm}

\LARGE{\textbf{Разработка и исследование RAG-системы для работы с PDF-документом}}
\vspace{2cm}
\end{center}
\begin{flushright}
\textbf{Выполнил:}\\
Студент группы МПИ-25-1-2\\
Г.В. Шевченко\\[0.5cm]
\textbf{Проверил:}\\
Преподаватель\\
М.Е. Корнет\\[1.5cm]
\textbf{Оценка:} \hrulefill
\end{flushright}
\vfill
\begin{center}Москва, 2026\end{center}
\end{titlepage}
\setcounter{page}{1}
\tableofcontents
\newpage
\section{Цель работы}
Цель работы --- построить воспроизводимый прототип поиска ответов по PDF на основе Retrieval Augmented Generation (RAG), исследовать влияние полноты индекса, размера и перекрытия фрагментов, сравнить два способа поиска, проверить ссылки на первоисточник, поведение при недостатке данных и поддержку уточнений в диалоге.

\section{Архитектура RAG-системы}
Исходный документ --- 12-страничное методическое пособие «Практическая работа №3: RAG-система». Для экспериментов индексировались страницы PDF 3--12 (10 страниц с основным материалом). Данные с ответами распределены по разным страницам. Исходное разбиение: 1000 символов, перекрытие 200 символов; индекс содержит 29 фрагментов, $k=3$.

\begin{figure}[H]\centering
\begin{tikzpicture}[node distance=1.2cm, every node/.style={font=\small}, box/.style={draw,rounded corners,align=center,minimum height=0.85cm,minimum width=3.0cm}, >=Stealth]
\node[box] (pdf) {PDF-документ\\PyMuPDF};
\node[box, right=of pdf] (split) {Разбиение\\метаданные};
\node[box, right=of split] (index) {TF--IDF + SVD\\или binary BoW};
\node[box, below=1.3cm of index] (search) {Поиск top-3\\косинусное сходство};
\node[box, left=of search] (context) {Контекст + ссылки\\страница и ID};
\node[box, left=of context] (ans) {Ответ + проверка\\опционально Ollama};
\draw[->] (pdf)--(split); \draw[->] (split)--(index);
\draw[->] (index)--(search); \draw[->] (search)--(context); \draw[->] (context)--(ans);
\node[above=0.12cm of search, xshift=-2.2cm] {Вопрос};
\end{tikzpicture}
\caption{Компонентная схема исследованного конвейера RAG}
\end{figure}

В зафиксированном \textbf{офлайн-прогоне} используется плотное представление на основе символьных TF--IDF-признаков и TruncatedSVD. Это \emph{не} нейросетевой \texttt{OllamaEmbeddings}. Вместо LLM применяется детерминированный извлекатель опорного фрагмента. Для запуска \textbf{настоящей LLM} предусмотрен отдельный вызов Ollama (модель \texttt{qwen2.5:3b}, температура 0); показанные ниже численные результаты не выдаются за генерацию LLM.

\begin{table}[H]\centering\caption{Файлы проекта}\small
\begin{tabular}{|L{5.0cm}|L{9.0cm}|}\hline
Файл & Назначение \\ \hline
\texttt{rag\_lab3.py} & Подготовка PDF, 7 экспериментов, индексы, подсчёт метрик и CSV \\ \hline
\texttt{Лабораторная\_работа\_3\_RAG.ipynb} & Notebook с сохранёнными выходными результатами \\ \hline
\texttt{run\_ollama\_demo.py} & Демонстрация одного реального ответа Ollama \\ \hline
\texttt{run\_ollama\_experiments.py} & Реальные ответы двух промптов, цитаты, уточнения; JSON/CSV \\ \hline
\texttt{test\_rag\_lab3.py} & Восемь автоматических тестов \\ \hline
\texttt{results\_lab3/} & Результаты задания 1--7 и сводка экспериментов \\ \hline
\end{tabular}
\end{table}

\section{Исходный код и запуск}
Проект полностью включён в сопровождающий архив. Запуск выполняется в каталоге с исходным PDF:
\begin{lstlisting}[language={},caption={Установка и проверка проекта}]
python -m pip install -r requirements_lab3.txt
python rag_lab3.py
python -m unittest -v test_rag_lab3.py
\end{lstlisting}
При установленном Ollama реальный ответ можно получить командами \texttt{ollama pull qwen2.5:3b} и \texttt{python run\_ollama\_demo.py}. Режимы экспериментов с LLM доступны в \texttt{run\_ollama\_experiments.py}. Подробности приведены в \texttt{README\_запуск.md}.

\section{Реализация заданий}
\subsection{Задание 1. Полный и сокращённый индекс}
Сопоставлены индексы из первых пяти и всех 29 фрагментов при неизменных вопросах и $k=3$. Банк содержит 10 вопросов: три по начальному разделу и семь по следующим. В каждой записи CSV сохранены ID найденных фрагментов и текст подтверждения при наличии.
\begin{table}[H]\centering\caption{Полнота поискового индекса}\begin{tabular}{|l|c|c|c|}\hline
Индекс & Начало, из 3 & Далее, из 7 & Hit Rate@3 \\ \hline
Первые 5 & 2 & 0 & 20\% \\ \hline
Полный & 1 & 6 & 70\% \\ \hline
\end{tabular}\end{table}
Два примера ошибки неполного индекса: вопросы Q06 (метаданные PDF-фрагмента) и Q08 (формула Hit Rate@k). Соответствующие опорные страницы лежат за пределами сокращённого набора; изменение промпта не создаёт отсутствующих в индексе сведений. Различие успеха на начальных вопросах показывает, что размер индекса также меняет ранжирование.

\subsection{Задание 2. Размер фрагментов}
Полный корпус был разбит при \texttt{chunk\_size}=500, 1000, 2000 символов; во всех случаях перекрытие 10 символов, top-3. Для каждой конфигурации индекс сформирован заново, вопросы остались прежними.
\begin{table}[H]\centering\caption{Влияние длины фрагментов}\begin{tabular}{|r|r|r|r|}\hline
Размер, символы & Фрагментов & Успехов из 10 & Hit Rate@3 \\ \hline
500 & 46 & 7 & 70\% \\ \hline
1000 & 25 & 6 & 60\% \\ \hline
2000 & 16 & 7 & 70\% \\ \hline
\end{tabular}\end{table}
Для этого набора вопросов результаты на 500 и 2000 символах совпали по агрегированной метрике, поэтому универсального победителя нет. Короткие фрагменты могут оборвать описание, длинные --- добавлять соседние темы. Результат на 10 вопросах не переносится автоматически на другие документы.

\subsection{Задание 3. Перекрытие фрагментов}
При \texttt{chunk\_size}=1000 и $k=3$ исследованы перекрытия 0, 100, 200 символов на пяти контрольных вопросах. Почти одинаковыми считались два извлечённых фрагмента с коэффициентом Жаккара множества символьных 5-грамм более 0,8.
\begin{table}[H]\centering\caption{Влияние перекрытия}\begin{tabular}{|r|r|r|r|}\hline
Перекрытие & Фрагментов & Успехов из 5 & Повторы в top-3 \\ \hline
0 & 25 & 4 & 0 \\ \hline
100 & 27 & 4 & 0 \\ \hline
200 & 29 & 5 & 0 \\ \hline
\end{tabular}\end{table}
Перекрытие 200 сохранило опорную цитату для дополнительного вопроса, но увеличило индекс. Улучшение не обязано быть монотонным: повторение текста конкурирует за места в top-3. Разбиение выполнялось отдельно внутри каждой страницы PDF.

\subsection{Задание 4. Плотный поиск и бинарный мешок слов}
Для одного набора фрагментов проверены пять вопросов с терминами из текста и пять переформулировок. В бинарном BoW использовался словарь, построенный только по PDF; запрос с нулевым вектором возвращает пустой результат. Плотная альтернатива получена методом SVD, а не отдельной моделью эмбеддингов.
\begin{table}[H]\centering\caption{Сравнение поисковых представлений}\begin{tabular}{|L{5.2cm}|c|c|c|}\hline
Метод & Точные из 5 & Иные формулировки из 5 & Всего из 10 \\ \hline
Бинарный мешок слов & 2 & 1 & 3 \\ \hline
TF--IDF + SVD & 3 & 3 & 6 \\ \hline
\end{tabular}\end{table}
В наблюдаемом офлайн-опыте SVD-поиск оказался успешнее по выбранной проверке опорных цитат. Это сравнение \emph{не подтверждает} превосходство нейросетевых эмбеддингов Ollama: они не участвовали в зафиксированном прогоне.

\subsection{Задание 5. Проверяемые ссылки на источники}
Каждый фрагмент содержит \texttt{id}, \texttt{page}, \texttt{source} и \texttt{start\_index}; позиция отсчитывается внутри извлечённого текста соответствующей страницы. Ссылка имеет формат [Фрагмент N, стр. P]. Приведены пять проверочных вопросов.
\begin{table}[H]\centering\caption{Проверка ссылок}\begin{tabular}{|c|c|c|c|}\hline
Вопрос & Ссылка существует & Факт подтверждён & Страница верна \\ \hline
Q01 & Нет & Нет & Нет \\ \hline
Q02 & Нет & Нет & Нет \\ \hline
Q03 & Да & Да & Да \\ \hline
Q04 & Да & Да & Да \\ \hline
Q05 & Да & Да & Да \\ \hline
\end{tabular}\end{table}
Итого 3 из 5 ответов снабжены проверенной ссылкой. Если top-3 не содержит точной опорной цитаты, извлекатель возвращает отказ, а не выдуманную ссылку. Формальная ссылка сама по себе не доказывает истинности утверждения: требуется проверять фрагмент и страницу.

\subsection{Задание 6. Отсутствие ответа и строгость промпта}
Сформированы пять вопросов, ответы на которые присутствуют в PDF, и пять тематически близких вопросов без ответа. Для обеих инструкций сохранены одинаковые найденные фрагменты. Исходный промпт: «Ответь на вопрос на основе найденных фрагментов. Используй ссылки, если возможно». Строгий промпт: «Отвечай исключительно по найденным фрагментам PDF. При недостатке информации явно откажись от ответа. Подтверждай факты ссылками».

\begin{table}[H]\centering\caption{Офлайн-прокси проверки ответа (не прогон двух LLM-промптов)}\begin{tabular}{|L{3cm}|c|c|c|}\hline
Промпт & Средний балл (5) & Отказы из 5 & Выдуманные ответы \\ \hline
Исходный & 1,2 & 5 & 0 \\ \hline
Строгий & 1,2 & 5 & 0 \\ \hline
\end{tabular}\end{table}
Оба результата одинаковы, поскольку использован \emph{один и тот же детерминированный извлекатель}, не реагирующий на изменение промпта. Следовательно, это не свидетельство о качестве двух инструкций для LLM; воспроизводимый сценарий их реального сравнения включён отдельно.

\subsection{Задание 7. Уточняющие вопросы}
Реализованы три последовательных диалога по темам RAG, бинарного мешка слов и памяти диалога. В каждом есть исходный вопрос и два уточнения, отдельный идентификатор сессии, хранение истории и очистка по команде. Переформулировка уточнений в самостоятельные запросы задана прозрачными правилами (вариант Б без LLM-переформулировщика).
\begin{table}[H]\centering\caption{Выполнение уточняющих вопросов}\begin{tabular}{|c|L{7.5cm}|c|}\hline
Сессия & Пример поискового запроса & Успехов из 2 \\ \hline
1 & Как работает подход RAG с поиском по документу? & 2 \\ \hline
2 & Учитывает ли бинарный мешок слов порядок слов? & 0 \\ \hline
3 & Сколько обменов вопрос--ответ хранит память диалога? & 2 \\ \hline
\end{tabular}\end{table}
Всего найдены опорные сведения в 4 из 6 уточнений. Проверено: истории трёх сессий независимы, очистка первой не влияет на вторую. История помогает восстановить объект вопроса, но не считается источником фактов.

\section{Тестирование}
\begin{table}[H]\centering\caption{Автоматические тест-кейсы}\begin{tabular}{|r|L{10cm}|l|}\hline
№ & Проверяемый сценарий & Итог \\ \hline
1 & Извлечение страниц PDF & PASS \\ \hline
2 & Создание более пяти фрагментов & PASS \\ \hline
3 & Корректные метаданные фрагмента & PASS \\ \hline
4 & Смещение при перекрытии 200 символов & PASS \\ \hline
5 & Ограничение выдачи top-3 & PASS \\ \hline
6 & Пустая выдача для вне-словарного запроса BoW & PASS \\ \hline
7 & Проверяемая ссылка по опорному фрагменту & PASS \\ \hline
8 & Генерация таблиц для семи заданий & PASS \\ \hline
\end{tabular}\end{table}
Локальный запуск \texttt{python -m unittest -q test\_rag\_lab3} завершился успешно: 8 тестов, ошибок нет. Запись результатов сохраняет метаданные и ответ (либо отсутствие подтверждения) для дополнительной ручной проверки.

\section{Вывод}
Реализован воспроизводимый офлайн-конвейер анализа PDF: загрузка, фрагментация, два вида представления, ранжирование top-3, ссылки, отказ при недостатке подтверждений и отдельные диалоговые сессии. Полный индекс повысил Hit Rate@3 с 20\% до 70\%, а перекрытие 200 символов дало 5 из 5 успешных поисков в эксперименте №3. На вопросах эксперимента №2 равный Hit Rate@3 показали фрагменты 500 и 2000 символов.

Ограничения существенны: поиск TF--IDF+SVD не равен OllamaEmbeddings; автоматическая проверка по опорной цитате не заменяет содержательную оценку генерации; влияние двух промптов на LLM в офлайн-прогоне не измерялось. Для самостоятельного запуска настоящей LLM и сохранения её реальных ответов приложены отдельные исполняемые скрипты.
\end{document}
